\section*{PRÉAMBULE}\label{pruxe9ambule}
\markboth{Préambule}{Préambule}
\addcontentsline{toc}{tocmain}{PRÉAMBULE}

\phantomsection\addcontentsline{toc}{section}{APRS Working Group}
\begin{aprssideblock}{APRS Working\\Group}

L'APRS Working Group est une association non constituée en personne
morale dont les membres s'engagent à faire progresser l'usage et à
accroître la valeur des protocoles APRS en : (a) publiant et maintenant
une spécification formelle du protocole APRS ; (b) publiant des tests de
validation et d'autres outils permettant la conformité à la
spécification ; (c) soutenant un programme de certification APRS ; (d)
travaillant plus généralement à améliorer les capacités d'APRS au sein
de la communauté radio-amateur.

Bien que le Working Group puisse recevoir un soutien du TAPR et d'autres
organisations, il s'agit d'un corps indépendant, non affilié à aucune
organisation. Le groupe n'a pas de budget, ne collecte pas de
cotisations et ne possède aucun actif.

Membres actuels de l'APRS Working Group :

\begin{tabular}{@{}lll@{}}
John Ackermann, & N8UR & Administrative Chair \& représentant TAPR \\
Bob Bruninga, & WB4APR & Technical Chair, fondateur d'APRS \\
Brent Hildebrand, & KH2Z & Auteur de APRS+SA \\
Stan Horzepa, & WA1LOU & Secrétaire \\
Mike Musick, & N0QBF & Auteur de pocketAPRS \\
Keith Sproul, & WU2Z & Co-auteur de WinAPRS/MacAPRS/X-APRS \\
Mark Sproul, & KB2ICI & Co-auteur de WinAPRS/MacAPRS/X-APRS \\
\end{tabular}
\end{aprssideblock}

\phantomsection\addcontentsline{toc}{section}{Remerciements}
\begin{aprssideblock}{Remerciements}

Ce document est le résultat de contributions de nombreuses personnes. Il
reprend une grande partie du matériel produit par des membres
individuels du Working Group.

De plus, le document sur le format Mic-E par Alan Crosswell N2YGK et Ron
Parsons W5RKN a servi de point de départ utile pour expliquer les
complications de ce format.
\end{aprssideblock}

\phantomsection\addcontentsline{toc}{section}{Numéro de version du document}
\begin{aprssideblock}{Numéro de\\version du\\document}

À l'exception de la toute première version publique en brouillon de la
référence du protocole APRS, le numéro de version du document est un
numéro à 3 parties \texttt{P.p.D} (pour une version approuvée) ou à 4
parties \texttt{P.p.Dd} (pour un brouillon) :

\begin{itemize}
\tightlist
\item
  \texttt{P} : version majeure du protocole APRS
\item
  \texttt{p} : version mineure du protocole APRS
\item
  \texttt{D} : release du document
\item
  \texttt{d} : lettre du brouillon
\end{itemize}

Ainsi, par exemple :

\begin{itemize}
\tightlist
\item
  « 1.2.3 » = release 3 du document couvrant la version 1.2 du protocole
  APRS.
\item
  « 1.2.3c » = brouillon « c » de ce document.
\end{itemize}
\end{aprssideblock}

\phantomsection\addcontentsline{toc}{section}{Historique des versions}
\begin{aprssideblock}{Historique des\\versions}

Voir Annexe 7.
\end{aprssideblock}

\phantomsection\addcontentsline{toc}{section}{Conventions du document}
\begin{aprssideblock}{Conventions\\du document}

\begin{itemize}
\tightlist
\item
  Police \texttt{Courier} : caractères ASCII dans les données APRS.
\item
  \texttt{V} : caractère espace ASCII.
\item
  \texttt{…} (ellipse) : zéro ou plusieurs caractères.
\item
  \texttt{/\$} : symbole de la table de symboles primaire.
\item
  \texttt{\textbackslash{}\$} : symbole de la table de symboles
  alternative.
\item
  \texttt{0x} : hexadécimal (ex. \texttt{0x1d}).
\item
  Tous les indicatifs sont supposés avoir un SSID \texttt{-0} sauf
  mention contraire.
\item
  \lit{Marqueur jaune} (fond gris clair en impression papier) : signale un
  texte d'intérêt --- particulièrement utile pour mettre en évidence des
  caractères ASCII littéraux uniques (par exemple \lit{"}) tels
  qu'ils apparaissent dans les données APRS.
\item
  Zones grisées dans les diagrammes de format de paquet : champs
  optionnels.
\end{itemize}
\end{aprssideblock}

\phantomsection\addcontentsline{toc}{section}{Retour}
\begin{aprssideblock}{Retour}

Merci d'adresser retours et commentaires sur ce document à la liste de
diffusion TAPR \texttt{aprsspec}. Pour s'y inscrire, partir de
\url{http://www.tapr.org} puis suivre le chemin \emph{Special Interest
Groups → APRS Specification → Join APRS Spec Discussion List}.
\end{aprssideblock}

\section*{AVANT-PROPOS DES AUTEURS}\label{avant-propos-des-auteurs}
\markboth{Avant-propos des auteurs}{Avant-propos des auteurs}
\addcontentsline{toc}{tocmain}{AVANT-PROPOS DES AUTEURS}

\begin{aprsbodyblock}
Ce document de référence décrit ce que l'on appelle la Version 1.0 du
protocole APRS, et constitue essentiellement une description du
fonctionnement actuel d'APRS.

Il est destiné principalement au programmeur souhaitant développer des
applications compatibles APRS, mais intéressera aussi l'utilisateur
ordinaire qui veut en savoir plus sur ce qui se passe « sous le capot ».

Il n'est cependant pas destiné à être une spécification de programmation
ennuyeuse, pédante, de style RFC, à lire et à comprendre uniquement par
les Mr Spocks de ce monde. Nous avons inclus de nombreux éléments
d'information générale qui, bien que ne faisant pas strictement partie
de la description formelle du protocole, fournissent un contexte utile
sur la façon dont APRS est réellement utilisé sur les ondes, et sur la
manière dont il est implémenté dans les logiciels APRS. Nous espérons
que cela mettra APRS en perspective, rendra le document plus lisible, et
n'offensera pas trop les puristes.

Il est important de comprendre comment APRS est né, et de saisir la
philosophie de conception derrière lui. En particulier, nous croyons
fermement qu'APRS est, et doit rester, un système tactique léger ---
presque n'importe qui devrait pouvoir l'utiliser dans des situations
temporaires (urgences, travail mobile, veille météo) avec un minimum de
formation et d'équipement.

Ce document est le résultat des contributions de nombreuses personnes,
collationnées et retravaillées par l'APRS Working Group. Nos
remerciements sincères à tous ceux qui ont contribué à sa mise en forme
et à sa publication. Si vous découvrez des erreurs, des omissions ou des
affirmations trompeuses, faites-le nous savoir --- le meilleur moyen est
la liste de diffusion TAPR \texttt{aprsspec} sur
\href{http://www.tapr.org}{www.tapr.org}.

Enfin, des utilisateurs du monde entier proposent en permanence de
nouvelles idées et suggestions pour étendre et améliorer APRS. Nous les
accueillons volontiers. Là encore, le meilleur endroit pour en discuter
est la liste \texttt{aprsspec}.

\emph{L'APRS Working Group}

\emph{Août 2000}
\end{aprsbodyblock}

\phantomsection\addcontentsline{toc}{section}{Avertissement}
\begin{aprssideblock}{Avertissement}

Comme tout système de navigation, APRS n'est pas infaillible. Nul ne
devrait s'appuyer aveuglément sur APRS pour naviguer, ou dans des
situations de vie ou de mort. De même, cette spécification n'est pas
infaillible.

Les membres de l'APRS Working Group ont fait de leur mieux pour définir
le protocole APRS, mais cette description peut contenir des erreurs ou
des omissions. Il est très probable que toutes les implémentations APRS
n'appliqueront pas pleinement ou correctement cette spécification,
aujourd'hui ou à l'avenir.

Nous exhortons quiconque utilise ou écrit un programme implémentant
cette spécification à faire preuve de prudence et de bon sens. L'APRS
Working Group et l'éditeur de la spécification déclinent toute
responsabilité pour tout dommage aux personnes ou aux biens qui pourrait
résulter de l'utilisation de cette spécification ou d'un logiciel
l'implémentant.
\end{aprssideblock}

\clearpage

\section*{STRUCTURE DE CETTE
SPÉCIFICATION}\label{structure-de-cette-spuxe9cification}
\markboth{STRUCTURE DE CETTE SPÉCIFICATION}{STRUCTURE DE CETTE SPÉCIFICATION}
\addcontentsline{toc}{tocmain}{STRUCTURE DE CETTE SPÉCIFICATION}

\begin{aprsbodyblock}
Cette spécification décrit les exigences globales pour développer un
logiciel conforme à la Version 1.0 du protocole APRS. Le flux
d'information part de la trame AX.25 UI standard et descend
progressivement dans les détails, l'usage de chaque champ de la trame
étant exploré.

Une caractéristique clé de la spécification est l'inclusion de dizaines
d'exemples détaillés de paquets APRS typiques et des calculs
mathématiques associés.

Voici le plan des chapitres :

\begin{list}{}{\setlength{\leftmargin}{0pt}\setlength{\rightmargin}{0pt}\setlength{\itemindent}{0pt}\setlength{\labelwidth}{0pt}\setlength{\labelsep}{0pt}\setlength{\itemsep}{0.35em}\setlength{\parsep}{0pt}}
\item
  \textbf{Introduction à APRS} --- bref historique et résumé des
  principales caractéristiques.
\item
  \textbf{La philosophie de conception APRS} --- les fondamentaux
  d'APRS, avec son usage comme outil de communication tactique temps
  réel, le timing des transmissions APRS et l'usage du digipeating
  générique.
\item
  \textbf{APRS et AX.25} --- bref rappel sur la structure de la trame
  AX.25 UI, avec référence particulière aux usages spéciaux qu'APRS fait
  des champs Destination, Source et Information.
\item
  \textbf{Données APRS dans les champs Destination et Source AX.25} ---
  détails des indicatifs APRS génériques et des indicatifs qui
  spécifient des symboles d'affichage et des numéros de version. Résumé
  du stockage des données Mic-E dans le champ Destination, et de la
  spécification d'une icône via le SSID Source.
\item
  \textbf{Données APRS dans le champ Information AX.25} --- détails des
  principaux constituants des données APRS stockées dans le champ
  Information. Contient la table des Data Type Identifiers APRS et un
  résumé de tous les types de données possibles.
\item
  \textbf{Formats de temps et de position} --- timestamps, latitude,
  longitude, ambiguïté de position, Maidenhead, NMEA, altitude.
\item
  \textbf{Extensions de données APRS} --- extensions optionnelles pour
  course/vitesse de la station, vitesse/direction du vent,
  puissance/hauteur/gain, portée radio pré-calculée, force de signal DF,
  descripteur d'objet zone.
\item
  \textbf{Formats de rapports de position et DF} --- détails complets de
  ces formats.
\item
  \textbf{Formats compressés de rapport de position} --- comment la
  position d'une station et les extensions APRS sont compressées en
  paquets très courts.
\item
  \textbf{Format de données Mic-E} --- encodage Mic-E de la lat/long,
  altitude, course, vitesse, code message Mic-E, données de télémétrie
  et chemin de digipeater APRS dans les champs Destination et
  Information de la trame AX.25.
\item
  \textbf{Rapports d'objets et d'éléments} --- comment configurer des
  objets et éléments APRS, et détails de l'encodage des Area Objects
  (cercles, lignes, ellipses, etc.).
\item
  \textbf{Rapports météo} --- détails complets pour les stations météo
  autonomes (sans position) et pour les rapports contenant une
  information de position. Détail du format Storm Data.
\item
  \textbf{Données de télémétrie} --- description du format MIM/KPC-3+,
  avec informations pour adapter l'interprétation des données brutes aux
  circonstances individuelles.
\item
  \textbf{Messages, bulletins et annonces} --- format complet.
\item
  \textbf{Capacités, requêtes et réponses de station} --- dix types de
  requêtes et les réponses attendues.
\item
  \textbf{Rapports de statut} --- format des messages de statut général,
  plus les cas particuliers d'usage pour véhiculer beam
  heading/puissance en meteor scatter et locator Maidenhead.
\item
  \textbf{Tunneling réseau} --- usage du Source Path Header pour
  tunneliser des paquets APRS à travers des réseaux tiers qui ne
  comprennent pas les adresses AX.25, et usage du Data Type Identifier
  tiers.
\item
  \textbf{Format défini par l'utilisateur} --- APRS permet aux
  utilisateurs de définir leurs propres formats à des fins
  particulières. Ce chapitre explique comment faire.
\item
  \textbf{Autres paquets} --- énoncé général sur la façon dont APRS doit
  traiter tout autre type de paquet non couvert par cette spécification.
\item
  \textbf{Symboles APRS} --- comment spécifier symboles et overlays dans
  les rapports de position et dans les indicatifs de destination GPS
  génériques.
\item
  \textbf{Formats de données APRS} --- annexe rassemblant tous les
  formats pour référence.
\item
  \textbf{Tables de symboles APRS} --- liste complète des symboles des
  tables Primary et Alternate.
\item
  \textbf{Table des codes ASCII} --- code ASCII complet, décimal et
  hexadécimal (le décimal est nécessaire pour les calculs de lat/long
  compressées et d'altitude), avec les codes hex des caractères ASCII
  bit-shiftés dans les adresses AX.25 (utile pour décoder Mic-E et pour
  le monitoring de paquets on-air).
\item
  \textbf{Glossaire} --- référence tout-en-un des nombreux termes
  spécifiques APRS.
\item
  \textbf{Références} --- pointeurs vers d'autres documents pertinents.
\end{list}
\end{aprsbodyblock}

\begin{center}\rule{0.5\linewidth}{0.5pt}\end{center}
